\SourceLeaf{282}
\Article{6}
Antequam quae hinc sequantur generaliter perscrutemur, ut modus argumentationis facilius penetrari possit, casus quosdam singulares evolvere visum est. Resumamus itaque primo formam $xx+yy$, pro qua itaque multitudo repraesentationum numeri indefiniti valorem medium $=\pi$ nanciscitur. Multitudo vero repraesentationum actualium numeri dati haud difficile e principiis generalibus in Disquisitionibus Arithmeticis stabilitis determinatur. Designemus per $fA$ multitudinem repraesentationum numeri $A$, quae erit $=4$, si $A=1$ vel $2$ vel potestas binarii; $=8$, si $A$ est numerus primus formae $4n+1$, vel productum
\SourceLeaf{283}
talis numeri primi in potestatem binarii; $=0$, si $A$ est numerus primus formae $4n+3$, vel per talem numerum primum divisibilis, neque vero per ipsius quadratum; denique \emph{generaliter}
\[
\begin{aligned}
\text{vel}\quad&=4(\alpha+1)(\beta+1)(\gamma+1)\ldots\\
\text{vel}\quad&=0
\end{aligned}
\]
prout, reducto numero $A$ ad formam $2^\mu Sa^\alpha b^\beta c^\gamma\ldots$, designantibus $a,b,c$ etc. numeros primos inaequales formae $4n+1$, $S$ autem productum e numeris primis formae $4n+3$, si qui inter factores numeri $A$ semel pluriesve occurrunt, numerus $S$ est vel quadratum vel non quadratum. Patet itaque, $fA$ unice pendere a modo, quo numeri primi $3,5,7,11,13$ etc. inter factores numeri $A$ reperiuntur, ita ut generaliter statuere oporteat
\[
fA=4(3)\pdot(5)\pdot(7)\pdot(11)\pdot(13)\dotsfour
\]
si valores characterum $(3),(5),(7)$ etc. ita acceptos supponimus, ut denotante $p$ numerum primum sit

\begin{description}
\item[primo] $(p)=1$, si $p$ ipsum $A$ non metitur
\item[secundo] $(p)=\alpha+1$, si $p$ est formae $4n+1$, atque $p^\alpha$ potestas summa ipsum $A$ metiens
\item[tertio] $(p)=0$, si $p$ est formae $4n+3$, atque exponens potestatis altissimae ipsius $p$ ipsum $A$ metientis est impar: denique
\item[quarto] $(p)=1$, si $p$ est formae $4n+3$, atque exponens potestatis summae ipsius $p$ ipsum $A$ metientis est par.
\end{description}
Manifesto casus primus sub secundo et quarto continetur.

Hoc itaque modo termini progressionis $f1,f2,f3,f4$ etc. valde irregulariter procedunt, etiamsi quo maior multitudo sumatur, eo accuratius valor medius $=\pi$ inde surgere debeat. Aggregatum $f1+f2+f3+\ldots+fA$ denotabimus per $FA$.

\Article{7}
Statuamus iam generaliter $fm+f3m=f'm$, perspicieturque facile, fieri
\[
f'A=4(5)\pdot(7)\pdot(11)\pdot(13)\ldots
\]
\SourceLeaf{284}
i. e. $f'A$ a relatione ipsius $A$ ad divisorem $3$ erit independens, unde seriei $f'1,f'2,f'3,f'4,f'5,f'6$ etc. irregularitas tum serius incipiet tum longe minor erit. Porro si statuimus
\[
f'1+f'2+f'3+f'4+\text{etc.}+f'm=F'm
\]
erit
\[
\begin{aligned}
F'3A&=F3A+f3+f6+f9+\ldots+f9A\\
&=F3A+FA
\end{aligned}
\]
Hinc facile concluditur crescente $A$ in infinitum, statui debere
\[
F'3A=4\pi A
\]
sive valorem medium terminorum seriei $f'1,f'2,f'3,f'4$ etc. esse
\[
=\tfrac43\pi
\]

Simili modo statuendo generaliter $-f'm+f'5m=f''m$, fiet
\[
f''A=4(7)(11)(13)\ldots
\]
sive e serie nova $f''1,f''2$ etc. abeunt vacillationes a relatione ad numerum $5$ pendentes. Statuendoque aggregatum
\[
f''1+f''2+f''3+\ldots+f''m=F''m
\]
fiet
\[
F''5m=-F'm+F'5m
\]
unde concluditur crescente $m$ in infinitum, statui debere
\[
F''5m=\tfrac43\pi\pdot4m
\]
sive valorem medium terminorum seriei esse $=\tfrac43\pdot\tfrac45\pi$.

Si eodem modo ulterius procedimus, progressiones novas formando, dum deinceps factores $(7),(11),(13),(17)$ etc. tollimus, hae continuo magis ad invariabilitatem appropinquabunt, valoresque medii deinceps novos factores $\tfrac87,\tfrac{12}{11},\tfrac{12}{13},\tfrac{16}{17},\tfrac{20}{19}$ etc. nanciscentur, ubi denominatores erunt numeri primi serie naturali, numeratores vero unitate vel maiores vel minores, prout illi sunt formae $4n-1$, vel ${}^{.}4n+1$. Quare quum hoc processu in infinitum continuato valor 
\SourceLeaf{285}
constans $4$ valori medio continuo propior fieri debeat, habemus
\[
4=\pi\pdot\tfrac43\pdot\tfrac45\pdot\tfrac87\pdot\tfrac{12}{11}\pdot\tfrac{12}{13}\ldots\text{ in inf.}
\]
sive
\[
\pi=4\pdot\frac3{3+1}\pdot\frac5{5-1}\pdot\frac7{7+1}\pdot\frac{11}{11+1}\pdot\frac{13}{13-1}\ldots
\]
Si singulae fractiones evolvuntur in series infinitas
\[
\begin{aligned}
\frac3{3+1}&=1-\frac13+\frac19-\frac1{27}+\ldots\\
\frac5{5-1}&=1+\frac15+\frac1{25}+\frac1{125}+\ldots\\
\frac7{7+1}&=1-\frac17+\frac1{49}-\frac1{343}+\ldots\\
&\text{etc.}
\end{aligned}
\]
productum facile evolvitur in
\[
1-\tfrac13+\tfrac15-\tfrac17+\tfrac19-\tfrac1{11}+\tfrac1{13}\dotsfour
\]
cuius seriei summam esse $=\tfrac14\pi$ vulgo notum est. Revera via inversa olim iam hinc aequalitas inter $\tfrac14\pi$ et productum infinitum $\tfrac43\pdot\tfrac45\pdot\tfrac87\ldotp\ldotp$ ab ill. \textsc{Euler} erutum fuerat (\emph{Introd. in analys. inf.} T. I. Cap. xv. art. 285).

\Article{8}
Consideremus secundo loco formam $xx+2yy$, pro qua multitudo repraesentationum numeri indefiniti valorem medium $=\frac\pi{\surd2}$ habebit. Designando per $fA$ multitudinem repraesentationum numeri dati $A$ per istam formam, haec erit $=2$ pro $A=1$ vel $A=2$, vel quoties $A$ est potestas binarii; porro $fA=4$, quoties $A$ est aliquis e serie numerorum primorum, quorum residuum quadraticum est $-2$, sive qui sunt formae $8k+1,8k+3$, puta $A=3,11,17,19,41,43$ etc.; denique $fA=0$, quoties $A$ est numerus primus, cuius non-residuum quadraticum est $-2$, puta e serie $5,7,13,23,29,31$ etc. sive vel formae $8k+5$, vel formae $8k+7$. Generaliter vero statui debet
\[
\begin{aligned}
\text{vel}\quad fA&=2(\alpha+1)(\beta+1)(\gamma+1)\ldots\\
\text{vel}\quad fA&=0
\end{aligned}
\]
prout reducto numero $A$ ad formam $2^\mu Sa^\alpha b^\beta c^\gamma\ldots$, designantibus $a,b,c$ etc. numeros primos inaequales formae $8k+1,8k+3$, contra $S$ productum e 
\SourceLeaf{286}
numeris reliquis (formae $8k+5,8k+7$), si qui inter factores numeri $A$ habentur, prout $S$ est quadratum vel non quadratum. Hinc per ratiocinia prorsus similia ut in art. praec. a serie $f1,f2,f3,f4,f5$ etc., puta $2,2,4,2,0,2$ etc. deinceps ad alias continuo longius \emph{constantes} progrediemur, quarum valores \emph{medii} sint deinceps $\frac\pi{\surd2}$, $\frac\pi{\surd2}\pdot\frac23$, $\frac\pi{\surd2}\pdot\frac23\pdot\frac65$, $\frac\pi{\surd2}\pdot\frac23\pdot\frac65\pdot\frac87$ etc.; progrediemur ita, ut deducamur ad aequationem
\[
2=\frac\pi{\surd2}\pdot\frac23\pdot\frac65\pdot\frac87\pdot\frac{10}{11}\pdot\frac{14}{13}\pdot\frac{16}{17}\pdot\frac{18}{19}\ldots
\]
ubi denominatores constituunt seriem naturalem numerorum primorum, numeratores vero unitate minores sunt, quoties denominatores sunt formae $8k+1$, vel $8k+3$, contra unitate maiores, quoties denominatores sunt formae $8k+5$ vel $8k+7$.
\BookRule{.30\textwidth}

